\section*{3. Data and variable construction}
\addcontentsline{toc}{section}{3. Data and variable construction}

This section documents the baseline and historical samples separately throughout. They are built from different DHS rounds, use different exposure constructions, and are validated independently; no claim in this section about one sample should be read as applying to the other. Numerical claims are sourced to the authoritative files listed in Appendix A; this section states what each variable is and how it was constructed, not the estimation results themselves (Section 5).

\subsection*{3.1 Baseline sample: unit of analysis, rounds, and pooling}
\addcontentsline{toc}{subsection}{3.1 Baseline sample: unit of analysis, rounds, and pooling}

The unit of analysis throughout the baseline sample is the child. The baseline pools the most recent DHS round available, as of data collection, for four countries: Ethiopia (2024--25), Ghana (2022), Kenya (2022), and Nigeria (2024). The full pooled, harmonised child file contains 70,231 children across 4,481 DHS clusters and 114 administrative-region fixed-effect cells (one set of regions per country, pooled). This pooled file is \textbf{not} the regression sample for any model: it is the universe before outcome-validity and exposure-classifiability restrictions are applied.

\subsection*{3.2 Baseline anthropometric outcomes and under-five eligibility}
\addcontentsline{toc}{subsection}{3.2 Baseline anthropometric outcomes and under-five eligibility}

Three outcomes are constructed following standard DHS/WHO anthropometric practice: height-for-age (HAZ), weight-for-age (WAZ), and weight-for-height (WHZ), each a standard-deviation z-score relative to the WHO Child Growth Standards reference population \parencite{deOnis2006}. A child contributes to an outcome only if that outcome's z-score is flagged valid in the DHS recode (implausible values, per DHS's own flagging convention, are excluded) and the child falls within the DHS-defined under-five anthropometric-measurement universe. Because valid measurements are not available for every child in the pooled file, the outcome-specific baseline samples are smaller than the full pooled file: 36,985 children (HAZ), 37,260 (WAZ), and 37,088 (WHZ). Comparing the household and community specifications on identical observations further requires both exposures to be classifiable for every child in the comparison, giving common baseline samples of 36,040 (HAZ), 36,308 (WAZ), and 36,145 (WHZ), spread across 4,426 DHS clusters. Under-five age is taken from the DHS-reported completed age in months where delivered directly by the recode (Ethiopia, Nigeria), and from a documented calendar-month difference construction otherwise (Ghana, Kenya), following each country round's own recode conventions.

\subsection*{3.3 Baseline WASH exposures: household and leave-one-out community}
\addcontentsline{toc}{subsection}{3.3 Baseline WASH exposures: household and leave-one-out community}

Four exposures are constructed for the baseline sample: own-household improved/non-improved water, own-household improved/non-improved sanitation, and their leave-one-out community counterparts, each following the WHO/UNICEF Joint Monitoring Programme (JMP) "improved" facility classification (piped water, protected wells/springs, and similar for water; sewer/septic and improved latrines for sanitation --- the same classification independently verified for the historical extension's gridded product in Section 3.4 and Appendix A). Household WASH is a household-level binary indicator. Community WASH is constructed for each child from the DHS household census of \textit{every other household in the same DHS cluster, excluding the child's own household} --- a leave-one-out construction, not an average over the other children in the analytical sample, so that two children in the same household never contribute to their own household's community measure and a child's own classification is never part of its own community exposure. This construction makes community WASH, by design, a property of the cluster that varies mainly \textit{between} clusters rather than within them; Section 4.3 derives precisely how little within-cluster variation the leave-one-out construction leaves and why this rules out a cluster-fixed-effects specification for the community exposure specifically, while it remains appropriate for the household exposure.

\subsection*{3.4 Baseline covariates, geography, and weights}
\addcontentsline{toc}{subsection}{3.4 Baseline covariates, geography, and weights}

Household- and child-level covariates include child sex, maternal age, maternal education (categorical), literacy, household size, birth order, wealth quintile, and urban/rural residence. DHS cluster (\texttt{IDHSPSU}, sample-qualified so that a cluster code never spans two survey rounds) and administrative region are retained as geographic identifiers; DHS cluster GPS coordinates are randomly displaced for confidentiality (urban clusters up to 2 km; rural clusters 99\% up to 5 km and 1\% up to 10 km; \parencite{Burgert2013}), which is not consequential for the baseline models (which do not use coordinates directly) but is directly consequential for the historical extension's spatial matching (Section 3.5). Baseline models use PERWEIGHT, DHS's standard child-level sampling weight, normalised within each model's own estimation sample so that each country contributes an equal total weight mass regardless of its raw sample size --- an authorised pooling choice, not a claim that any single country's raw sample size should determine its influence on a pooled coefficient (Section 4.5 gives the exact formula and its consequences).

\subsection*{3.5 Historical extension: rounds, exposure source, and scope}
\addcontentsline{toc}{subsection}{3.5 Historical extension: rounds, exposure source, and scope}

The historical extension links an earlier DHS round for three countries --- Ghana (2014), Kenya (2014), and Nigeria (2018) --- to IHME's annual gridded estimates of local access to improved drinking water (W\_IMP) and improved sanitation (S\_IMP), available for 2000--2017 \parencite{LBDWaSH2020}. These product definitions were verified directly from the primary source during this project's development, not inferred from the 0--100 raster scale or file names alone: W\_IMP is the sum of piped water (on- or off-premises) and other JMP-improved sources (protected wells and springs, bottled water, rainwater collection, bought water); S\_IMP is the sum of sewer/septic sanitation and other JMP-improved facilities (improved and ventilated improved latrines, composting toilets), each reported by the source as percent of population with access on a continuous \textasciitilde{}5×5 km surface. This project's own spatial extraction computes an \textbf{area average of the raster surface within each child's exposure buffer}, not the source study's own population-weighted administrative aggregate; this distinction is maintained throughout and restated wherever a historical coefficient is reported (Section 5.3; Section 6). Ethiopia's 2016 round is included only as a separately labelled, provisional sensitivity analysis (Appendix C), not as part of the established-date pooled sample, because the Ethiopian-calendar-to-Gregorian interview-date conversion required to place Ethiopian children's exposure windows could not be validated to the same standard as the three established-date rounds (Appendix C details the three candidate conversions and why one of them, labelled C1, is invalid and excluded from reporting). As stated in Section 1.2, this construction assigns each child a \textit{modelled, area-level} local coverage estimate for a defined early-life window; it is a measurement improvement on the timing dimension specifically, not a record of what that child's own household facilities actually were, which this or any other available data source cannot provide.

\subsection*{3.6 Historical spatial extraction: buffers, displacement, and coverage}
\addcontentsline{toc}{subsection}{3.6 Historical spatial extraction: buffers, displacement, and coverage}

Each child's exposure buffer is centred on its DHS cluster's (displaced) coordinates using an azimuthal equidistant projection local to that cluster, which avoids the area-distortion that a single fixed projection introduces at different latitudes. Within that buffer, the raster's value is averaged using an exact disk--pixel overlap computation (not a centroid or nearest-neighbour approximation), so that a pixel partially inside the buffer contributes its exact area-weighted share. Coverage completeness is checked explicitly: a child-year is assigned a valid extraction only when the buffer's valid-pixel area meets a documented tolerance (within 10⁻³ of the full buffer area), and partial- or missing-coverage child-years are excluded from that window rather than silently assigned a value computed from incomplete coverage.

\subsection*{3.7 Historical exposure windows and completion rules}
\addcontentsline{toc}{subsection}{3.7 Historical exposure windows and completion rules}

Four candidate exposure windows are constructed for each child, applied to the same underlying annual raster series: the calendar year of birth; the nine months before birth (prenatal); the first 12 completed postnatal months (the working main window; Section 4.1 states why); and the first 24 postnatal months. Because IHME's product is an annual surface, a window spanning two calendar years is assigned a weighted average of the relevant years' values, weighted by how many of the window's months fall in each calendar year. A child contributes to a given window only if every calendar year the window touches has a verified raster value \textbf{and} the window has been completed by the child's interview date (so a child interviewed before their own first-year postnatal window has elapsed is excluded from that window, not assigned a partial-year value). This completion rule means that, even for windows whose exposure \textit{value} does not depend on the precise interview date (Appendix C shows this is true for the birth-year window), eligibility for inclusion still requires a valid interview-date conversion --- relevant to why Ethiopia's three interview-date candidates matter even for windows that are otherwise interview-date-insensitive.

\subsection*{3.8 Historical covariates, weights, and the established-date measurement universe}
\addcontentsline{toc}{subsection}{3.8 Historical covariates, weights, and the established-date measurement universe}

Historical models control for cluster fixed effects, survey-qualified country-by-birth-year terms (each country's own birth-year effects relative to that country's own earliest observed birth year, not a single pooled reference --- Section 4.5 explains why this distinction matters for the pooled estimate), child sex, maternal age, maternal education, and marital status, with the child's own completed age in months (reported where delivered, documented calendar-month construction otherwise, matching the baseline's rule in Section 3.2) entered as an additional linear control. Household wealth quintile is \textbf{excluded from the historical extension's main specification}, for a documented, non-fixed-effects reason: current wealth may be a mediator or a post-exposure outcome of the historical WASH exposure being studied, and wealth quintile is correlated with the WASH exposure itself, so including it in the main specification risks conditioning on a variable downstream of the exposure of interest; wealth is included only as a labelled sensitivity check (Section 5.3), where the pooled water--HAZ coefficient is essentially unchanged (1.01 vs. 1.03). This is a different, and more defensible, reason than "cluster fixed effects absorb wealth" --- that claim is false in general, since cluster fixed effects absorb only what is constant \textit{within} a cluster, and wealth varies across households within the same cluster just as WASH status does (confirmed directly for this thesis's own baseline cluster-fixed-effects model, where wealth quintile \textit{is} included and estimated as a standard covariate; Section 4.4).

Weighting uses PERWEIGHT, normalised so that each of the three established-date countries contributes an equal total observation weight within each model's own final estimation sample (not population-size weighting) --- the same authorised convention as the baseline, carried over for comparability, and reported alongside its unweighted counterpart in every case (Section 5.3). The documented sampling weight for measured anthropometry specifically --- as distinct from the general household/individual sampling weight --- could not be located in official DHS documentation despite repeated attempts at alternative official sources (Appendix A records the exact sources tried and their outcomes); this project's working inference, stated explicitly wherever historical coefficients are reported, is that within-survey PERWEIGHT normalisation removes a constant measurement-subsample selection factor under an assumption of random household-level selection into the anthropometry subsample, but it cannot and does not correct for any differential non-response among eligible children who were not measured. Historical results are therefore reported as \textbf{weighted conditional associations among eligible measured children}, not as population-representative estimates for any defined population, and this phrasing is used consistently rather than varied across sections.

\subsection*{3.9 A documented measurement-universe discrepancy: Nigeria 2018}
\addcontentsline{toc}{subsection}{3.9 A documented measurement-universe discrepancy: Nigeria 2018}

Three distinct child populations arise for Nigeria 2018, and this thesis does not force them to a common count, nor treat any apparent gap between them as resolved. The official 2018 NDHS final report counts 12,806 (unweighted) children under five as eligible for measurement in the survey's anthropometry subsample (households selected for the men's questionnaire, roughly one-third of all households) and actually measured. The IPUMS-DHS birth-history extract used in this thesis contains 33,924 birth-history records for children born within approximately five years of the survey, of which 11,704 are measured or flagged non-missing for anthropometry. The final historical-extension estimation sample (water, HAZ, the main post-12-month window) retains 6,912 children after applying the window-completion, spatial-coverage, and singleton-cluster exclusions described above. These three counts differ by construction --- the official count is eligible children in subsample households; the extract is birth-history children of interviewed women, which can omit children whose mothers are not themselves in the sampled frame (a hypothesis raised in this project's own documentation, not independently verified); and the estimation sample further requires complete exposure windows and complete spatial coverage. The gap between the official count and the extract (12,806 vs. 11,704 non-missing) is disclosed here as an open, unresolved measurement-universe question and is not treated as resolved by any estimation choice in this thesis.
